% Generated draft: requires mathematical and source review.
\documentclass[11pt]{article}
\usepackage[margin=25mm]{geometry}
\usepackage{fontspec}
\setmainfont{Latin Modern Roman}
\usepackage{amsmath,amssymb,mathtools,unicode-math}
\setmathfont{Latin Modern Math}
\usepackage{enumitem,needspace,xurl,hyperref}
\setlength{\parindent}{0pt}
\setlength{\parskip}{5pt}
\setlength{\emergencystretch}{3em}
\newcommand{\N}{\mathbb{N}}
\newcommand{\Z}{\mathbb{Z}}
\newcommand{\Q}{\mathbb{Q}}
\newcommand{\R}{\mathbb{R}}
\newcommand{\C}{\mathbb{C}}
\newcommand{\F}{\mathbb{F}}
\newcommand{\A}{\mathbb{A}}
\newcommand{\PP}{\mathbb{P}}
\newcommand{\E}{\mathbb{E}}
\newcommand{\eps}{\varepsilon}
\newcommand{\dd}{\,\mathrm{d}}
\begin{document}
\section*{UnsolvedMath ID: 20002073}
\textbf{Source problem code:} catalogue entry
\subsection{Definitions and mathematical statement}
Question 1.15 has two parts: (a) develop singular ``Weinstein spaces'' analogous to Stein spaces, and (b) decide connectedness of the space of plurisubharmonic functions on a Weinstein manifold, with no critical points at infinity, also asking about fixed contact structure at infinity. The source additionally mentions an unspecified manifold diffeomorphic to \(\mathbb C^n\) for which every relevant plurisubharmonic function needs more than one critical point. Its exact function class and example are not supplied.

A plurisubharmonic condition needs a chosen complex or almost complex structure \(J\); a Weinstein manifold by itself does not fix all such data. In the aid's fixed-\(J\) smooth setting, a strict potential satisfies positivity of its Levi form, written \(dd^c\phi(v,Jv)>0\) in a consistent sign convention, and exhausting means proper and bounded below. The aid reports that the nonempty space of smooth exhausting strict \(J\)-plurisubharmonic functions is convex and contractible. The subspace agreeing with a prescribed exterior germ outside one fixed compact set is also convex. If that germ has no critical points there, interpolation preserves that particular no-critical-points-at-infinity condition. Such a condition without fixed exterior control is not asserted automatically preserved by convex combinations.

For a different fixed-data problem, the aid fixes a Liouville vector field \(Z\) and local Morse germs at its zeros. Exhausting Lyapunov functions with those germs and
\[
 d\phi(Z)>0\qquad\text{away from the zeros}
\]
again form a convex class under its stated conditions. This does not claim contractibility when all Weinstein or Stein structures are allowed to vary simultaneously.

The aid contrasts a Morse-only class even on \(\mathbb C\) with fixed standard \(J\). It reports a path of exhausting strict potentials agreeing with the same complete exterior potential outside one compact rectangle, with endpoints having respectively one minimum, and two minima plus one saddle. Along the displayed path a quartic Morse degeneracy occurs; a small suitable path perturbation can instead yield a Cerf-generic birth--death. The associated Liouville forms also agree outside that rectangle. The explicit potential and cutoff formulas are not reproduced in the supplied record.

Within this fixed-exterior class all critical points remain in a common compact set. Their number is locally constant along any all-Morse path, so the reported endpoints lie in different components of the all-Morse subspace even though the full strict-potential space is contractible. Allowing degenerate critical points removes this particular obstruction. Fixed exterior potential/form is stronger than fixing only contact structure at infinity; the restricted counterexample alone does not prove disconnectedness in every larger end-equivalence space where critical points might escape.

For part (a), the report describes arboreal skeleta, Liouville sectors and cotangent buildings as partial models, not a complete singular ambient counterpart of Stein spaces. It leaves that program and the original variable-data interpretation of part (b) unresolved.
\subsection{Short English statement}
The supplied aid separates fixed-complex-structure potentials, fixed-Liouville Lyapunov functions and Morse-only paths. The first two appropriate fixed-data classes are convex; a fixed-exterior Morse-only class on \(\mathbb C\) has a reported critical-count obstruction. This does not settle every meaning of the original Weinstein-space or end-fixed connectivity questions.
\subsection{Sources}
\begin{itemize}
\item Local source record: catalogue entry, catalogue entry; canonical \texttt{aim-geometry-notes.json}, \texttt{notes[410]}, Question 1.15. Title: ``Fixed-end connectivity and Morse walls for Weinstein potentials.''
\item Michael Hutchings, with help from Yasha Eliashberg and John Etnyre, \emph{Holomorphic curves in contact geometry}, AIM workshop problem list, version 15 October 2003; supplied reference \url{https://aimath.org/WWN/contactgeom2/contactgeom2.pdf}. This draft uses the supplied record without accessing a PDF.
\item Supplied research attempt 1: fixed-data convexity trichotomy, compactly supported potential deformation, quartic/Cerf distinction and partial singular-Weinstein models. No full external bibliographic entries, explicit potential formulas or identification of the source's exotic \(\mathbb C^n\) example are supplied.
\item Registry dated 14 August 2026: ``partially\_solved,'' ``PARTIAL-PROGRESS,'' statement status ``exact''; attempt partial, rigor ``proved,'' novelty confidence low. The machine-generated, non-peer-reviewed aid requires independent verification before citation.
\end{itemize}
\end{document}
